\makeatletter
\def\rap@fontpath{../../../Common/}
\makeatother
\documentclass{../../../Common/Latex/Templates/rapport}
% ---------------------------------------------------------------------------
% Shared settings for the H2026 General Assembly document set.
%
% Usage, directly after \documentclass{../../../Common/Latex/Templates/rapport}
% in a document one folder below this file:
%   % ---------------------------------------------------------------------------
% Shared settings for the H2026 General Assembly document set.
%
% Usage, directly after \documentclass{../../../Common/Latex/Templates/rapport}
% in a document one folder below this file:
%   % ---------------------------------------------------------------------------
% Shared settings for the H2026 General Assembly document set.
%
% Usage, directly after \documentclass{../../../Common/Latex/Templates/rapport}
% in a document one folder below this file:
%   \input{../ga-common}
%
% Image paths are relative to the compiled document's folder, so every
% document using this file must sit one folder below it.
% ---------------------------------------------------------------------------

% Meeting details, used on covers and in each document's ID block.
\newcommand{\meetingdate}{30 September 2026}
\newcommand{\meetingtime}{12:00}
\newcommand{\meetinglocation}{Festsalen, NMBU}
\newcommand{\meetingline}{\textbf{Meeting:} \meetingdate, \meetingtime, \meetinglocation}

% Header: moon logo on the left, full-text logo on the right. \projectname is
% still set because the cover title uses it.
\projectname{Eclipse NMBU}
\projectlogo{../../../Common/Eclipse Images/EclipseMoon}
\orglogo{../../../Common/Eclipse Images/EclipseFullText}

%
% Image paths are relative to the compiled document's folder, so every
% document using this file must sit one folder below it.
% ---------------------------------------------------------------------------

% Meeting details, used on covers and in each document's ID block.
\newcommand{\meetingdate}{30 September 2026}
\newcommand{\meetingtime}{12:00}
\newcommand{\meetinglocation}{Festsalen, NMBU}
\newcommand{\meetingline}{\textbf{Meeting:} \meetingdate, \meetingtime, \meetinglocation}

% Header: moon logo on the left, full-text logo on the right. \projectname is
% still set because the cover title uses it.
\projectname{Eclipse NMBU}
\projectlogo{../../../Common/Eclipse Images/EclipseMoon}
\orglogo{../../../Common/Eclipse Images/EclipseFullText}

%
% Image paths are relative to the compiled document's folder, so every
% document using this file must sit one folder below it.
% ---------------------------------------------------------------------------

% Meeting details, used on covers and in each document's ID block.
\newcommand{\meetingdate}{30 September 2026}
\newcommand{\meetingtime}{12:00}
\newcommand{\meetinglocation}{Festsalen, NMBU}
\newcommand{\meetingline}{\textbf{Meeting:} \meetingdate, \meetingtime, \meetinglocation}

% Header: moon logo on the left, full-text logo on the right. \projectname is
% still set because the cover title uses it.
\projectname{Eclipse NMBU}
\projectlogo{../../../Common/Eclipse Images/EclipseMoon}
\orglogo{../../../Common/Eclipse Images/EclipseFullText}

\usepackage{enumitem}
\usepackage{array}
\usepackage{tabularx}
\usepackage{amssymb}

\reporttitle{Responsibility \& Signature Contract}

% ---------------------------------------------------------------------------
% Cover / title page for rapport-class documents.
%
% Usage (after \documentclass{rapport} and before \begin{document}):
%   % ---------------------------------------------------------------------------
% Cover / title page for rapport-class documents.
%
% Usage (after \documentclass{rapport} and before \begin{document}):
%   % ---------------------------------------------------------------------------
% Cover / title page for rapport-class documents.
%
% Usage (after \documentclass{rapport} and before \begin{document}):
%   \input{../cover/rapport-cover}
%   \missionpatch{path/to/patch}      % optional circular mission/team patch;
%                                     % without it, the org logo is shown large
%                                     % at the top instead
%   \coverkind{Vedtekter}             % optional small line above the name
%   \reportsubtitle{Optional subtitle line}
%   \orgname{Eclipse NMBU}{Norwegian University of Life Sciences, Norway}
%   \reportlocation{Ås, Norge}
%   \date{3. oktober 2022}            % optional, defaults to \today
%   \contactinfo{Contact: post@eclipsenmbu.no}  % optional, shown at the bottom
%
% Then, as the first thing inside \begin{document}:
%   \makecoverpage
%
% The large title line shows \projectname (falling back to \reporttitle if
% \coverkind is not set), so a document typically sets both:
%   \projectname{Eclipse NMBU}
%   \coverkind{Vedtekter}
% ---------------------------------------------------------------------------

\makeatletter
\def\@missionpatch{}
\newcommand{\missionpatch}[1]{\def\@missionpatch{#1}}

\def\@coverkind{}
\newcommand{\coverkind}[1]{\def\@coverkind{#1}}

\def\@reportsubtitle{}
\newcommand{\reportsubtitle}[1]{\def\@reportsubtitle{#1}}

\def\@orgnameone{}
\def\@orgnametwo{}
\newcommand{\orgname}[2]{\def\@orgnameone{#1}\def\@orgnametwo{#2}}

\def\@reportlocation{}
\newcommand{\reportlocation}[1]{\def\@reportlocation{#1}}

\def\@contactinfo{}
\newcommand{\contactinfo}[1]{\def\@contactinfo{#1}}

\newcommand{\makecoverpage}{%
  \begingroup
  \thispagestyle{empty}
  \begin{center}
  \ifx\@missionpatch\@empty
    \includegraphics[width=10cm]{\@orglogo}\par
  \else
    \includegraphics[width=6.5cm]{\@missionpatch}\par
  \fi
  \vspace{3.4cm}
  \ifx\@coverkind\@empty
    {\sffamily\Huge\bfseries \@reporttitle \par}
  \else
    {\sffamily\Huge\bfseries \@coverkind \par}
    {\sffamily\Huge\bfseries \@projectname \par}
  \fi
  \ifx\@reportsubtitle\@empty\else
    \vspace{1.6cm}
    {\sffamily\bfseries \@reportsubtitle \par}
  \fi
  \vspace{2cm}
  \ifx\@missionpatch\@empty\else
    \includegraphics[height=1.3cm]{\@orglogo}\par
    \vspace{0.6cm}
  \fi
  \@orgnameone \\
  \@orgnametwo \par
  \vspace{0.6cm}
  \ifx\@reportlocation\@empty
    \@date
  \else
    \@reportlocation, \@date
  \fi
  \par
  \vfill
  \ifx\@contactinfo\@empty\else
    {\small \@contactinfo \par}
  \fi
  \end{center}
  \endgroup
  \newpage
  \pagenumbering{roman}
  \pagestyle{plain}
}
\makeatother

%   \missionpatch{path/to/patch}      % optional circular mission/team patch;
%                                     % without it, the org logo is shown large
%                                     % at the top instead
%   \coverkind{Vedtekter}             % optional small line above the name
%   \reportsubtitle{Optional subtitle line}
%   \orgname{Eclipse NMBU}{Norwegian University of Life Sciences, Norway}
%   \reportlocation{Ås, Norge}
%   \date{3. oktober 2022}            % optional, defaults to \today
%   \contactinfo{Contact: post@eclipsenmbu.no}  % optional, shown at the bottom
%
% Then, as the first thing inside \begin{document}:
%   \makecoverpage
%
% The large title line shows \projectname (falling back to \reporttitle if
% \coverkind is not set), so a document typically sets both:
%   \projectname{Eclipse NMBU}
%   \coverkind{Vedtekter}
% ---------------------------------------------------------------------------

\makeatletter
\def\@missionpatch{}
\newcommand{\missionpatch}[1]{\def\@missionpatch{#1}}

\def\@coverkind{}
\newcommand{\coverkind}[1]{\def\@coverkind{#1}}

\def\@reportsubtitle{}
\newcommand{\reportsubtitle}[1]{\def\@reportsubtitle{#1}}

\def\@orgnameone{}
\def\@orgnametwo{}
\newcommand{\orgname}[2]{\def\@orgnameone{#1}\def\@orgnametwo{#2}}

\def\@reportlocation{}
\newcommand{\reportlocation}[1]{\def\@reportlocation{#1}}

\def\@contactinfo{}
\newcommand{\contactinfo}[1]{\def\@contactinfo{#1}}

\newcommand{\makecoverpage}{%
  \begingroup
  \thispagestyle{empty}
  \begin{center}
  \ifx\@missionpatch\@empty
    \includegraphics[width=10cm]{\@orglogo}\par
  \else
    \includegraphics[width=6.5cm]{\@missionpatch}\par
  \fi
  \vspace{3.4cm}
  \ifx\@coverkind\@empty
    {\sffamily\Huge\bfseries \@reporttitle \par}
  \else
    {\sffamily\Huge\bfseries \@coverkind \par}
    {\sffamily\Huge\bfseries \@projectname \par}
  \fi
  \ifx\@reportsubtitle\@empty\else
    \vspace{1.6cm}
    {\sffamily\bfseries \@reportsubtitle \par}
  \fi
  \vspace{2cm}
  \ifx\@missionpatch\@empty\else
    \includegraphics[height=1.3cm]{\@orglogo}\par
    \vspace{0.6cm}
  \fi
  \@orgnameone \\
  \@orgnametwo \par
  \vspace{0.6cm}
  \ifx\@reportlocation\@empty
    \@date
  \else
    \@reportlocation, \@date
  \fi
  \par
  \vfill
  \ifx\@contactinfo\@empty\else
    {\small \@contactinfo \par}
  \fi
  \end{center}
  \endgroup
  \newpage
  \pagenumbering{roman}
  \pagestyle{plain}
}
\makeatother

%   \missionpatch{path/to/patch}      % optional circular mission/team patch;
%                                     % without it, the org logo is shown large
%                                     % at the top instead
%   \coverkind{Vedtekter}             % optional small line above the name
%   \reportsubtitle{Optional subtitle line}
%   \orgname{Eclipse NMBU}{Norwegian University of Life Sciences, Norway}
%   \reportlocation{Ås, Norge}
%   \date{3. oktober 2022}            % optional, defaults to \today
%   \contactinfo{Contact: post@eclipsenmbu.no}  % optional, shown at the bottom
%
% Then, as the first thing inside \begin{document}:
%   \makecoverpage
%
% The large title line shows \projectname (falling back to \reporttitle if
% \coverkind is not set), so a document typically sets both:
%   \projectname{Eclipse NMBU}
%   \coverkind{Vedtekter}
% ---------------------------------------------------------------------------

\makeatletter
\def\@missionpatch{}
\newcommand{\missionpatch}[1]{\def\@missionpatch{#1}}

\def\@coverkind{}
\newcommand{\coverkind}[1]{\def\@coverkind{#1}}

\def\@reportsubtitle{}
\newcommand{\reportsubtitle}[1]{\def\@reportsubtitle{#1}}

\def\@orgnameone{}
\def\@orgnametwo{}
\newcommand{\orgname}[2]{\def\@orgnameone{#1}\def\@orgnametwo{#2}}

\def\@reportlocation{}
\newcommand{\reportlocation}[1]{\def\@reportlocation{#1}}

\def\@contactinfo{}
\newcommand{\contactinfo}[1]{\def\@contactinfo{#1}}

\newcommand{\makecoverpage}{%
  \begingroup
  \thispagestyle{empty}
  \begin{center}
  \ifx\@missionpatch\@empty
    \includegraphics[width=10cm]{\@orglogo}\par
  \else
    \includegraphics[width=6.5cm]{\@missionpatch}\par
  \fi
  \vspace{3.4cm}
  \ifx\@coverkind\@empty
    {\sffamily\Huge\bfseries \@reporttitle \par}
  \else
    {\sffamily\Huge\bfseries \@coverkind \par}
    {\sffamily\Huge\bfseries \@projectname \par}
  \fi
  \ifx\@reportsubtitle\@empty\else
    \vspace{1.6cm}
    {\sffamily\bfseries \@reportsubtitle \par}
  \fi
  \vspace{2cm}
  \ifx\@missionpatch\@empty\else
    \includegraphics[height=1.3cm]{\@orglogo}\par
    \vspace{0.6cm}
  \fi
  \@orgnameone \\
  \@orgnametwo \par
  \vspace{0.6cm}
  \ifx\@reportlocation\@empty
    \@date
  \else
    \@reportlocation, \@date
  \fi
  \par
  \vfill
  \ifx\@contactinfo\@empty\else
    {\small \@contactinfo \par}
  \fi
  \end{center}
  \endgroup
  \newpage
  \pagenumbering{roman}
  \pagestyle{plain}
}
\makeatother

\coverkind{Responsibility \& Signature Contract (ansvars- og signaturkontrakt)}
\reportsubtitle{New Bylaws version -- 26H-ECL-GA-RESP1B}
\orgname{Eclipse NMBU}{Norwegian University of Life Sciences, Norway}
\reportlocation{\meetinglocation, Ås}
\date{\meetingdate, \meetingtime}
\contactinfo{Contact: post@eclipsenmbu.no}

\begin{document}

\makecoverpage
\tableofcontents
\reportmainmatter

\noindent\textbf{Document ID:} 26H-ECL-GA-RESP1B\\
\meetingline\\
\textbf{Applies from:} the New Bylaws' entry into force, on signature of the Meeting Protocol (26H-ECL-GA-PROT1). While the Old Bylaws govern, 26H-ECL-GA-RESP1A is used.\\
\textbf{Governing Bylaws:} proposed New Bylaws \S14 (Mandate, delegation and Signature Right), \S4.3 no.\ 7--8 (Nearest Leader, Responsibility Role), \S5.1 (Core Member).

\vsection{1. Core principle}
Signing on the Association's behalf is itself a responsibility, not just an administrative permission. Signature Right is assigned through this same contract, alongside any other genuine leadership responsibility: Responsibility, delegation, Authority and Signature Right are one relationship.

\vsection{2. Parties}
\begin{tabularx}{\textwidth}{|>{\bfseries\raggedright\arraybackslash}p{3cm}|X|}
  \hline
  From & \rule{0pt}{1cm}The person currently holding the responsibility (or, where the General Assembly itself is the granting party, see \S8 below). \\
  \hline
  To   & \rule{0pt}{1cm}The person receiving the responsibility.                                                                                      \\
  \hline
\end{tabularx}

Signed by both parties, then logged by the Board (Board Record, 26H-ECL-GA-RESP2). \textbf{Effective automatically on signature by both parties}: the Board's role is to record the transfer, not approve it. The Board keeps its ordinary powers under \S5.7 and \S14.5 (e.g. to change who holds an office, or to suspend a delegated Mandate by 2/3 vote) independently of this contract.

\vsection{3. Responsibility direction}
Responsibility moves only \textbf{downward} (From a leader To someone reporting to them, directly or via the chain in proposed New Bylaws \S4.3 no.\ 7), never sideways between peers. It moves upward only through Escalation (\S6), when the responsible superior is unavailable -- not as an ordinary transfer.

\vsection{4. Responsibility being transferred}
\begin{tabularx}{\textwidth}{|>{\bfseries\raggedright\arraybackslash}p{4.5cm}|X|}
  \hline
  Scope                                & \rule{0pt}{1.6cm}                                                                                                                                             \\
  \hline
  Specific responsibilities            & \rule{0pt}{1.8cm}                                                                                                                                             \\
  \hline
  Expected outcomes                    & \rule{0pt}{1.6cm}                                                                                                                                             \\
  \hline
  Delivery obligation (leveringsplikt) & \rule{0pt}{1.2cm}The recipient is obliged to deliver within the set time frame (proposed New Bylaws \S5.6.2).                                                 \\
  \hline
  Reporting obligation (meldeplikt)    & \rule{0pt}{1.2cm}If the recipient cannot carry out a task taken on, they shall notify their Nearest Leader without undue delay (proposed New Bylaws \S5.6.3). \\
  \hline
  Instance-specific upward obligations & \rule{0pt}{1.6cm}Any additional reporting or approval steps specific to this responsibility (e.g. financial limits, safety sign-off) -- list here.            \\
  \hline
  Signature Right included?            & \rule{0pt}{1cm}Yes \quad $\square$ \quad\quad No \quad $\square$ \quad\quad If yes: within what limits? \rule{3cm}{0.4pt} (proposed New Bylaws \S14.9)        \\
  \hline
\end{tabularx}

\vsection{5. Obligations of From}
The recipient (To) can expect from the person giving responsibility (From):
\begin{enumerate}
  \item Clear intent: what is being handed over, and why.
  \item Real delegated Authority, not just a task -- the standing to make the decisions the responsibility requires.
  \item The resources and access the responsibility requires.
  \item Availability for questions and support.
  \item Support where the recipient's own team falls short.
  \item Honest feedback on performance.
\end{enumerate}

\vsection{6. Escalation}
If \textbf{From} is unavailable, escalation goes to From's \textbf{Nearest Leader (nærmeste leder)}, per proposed New Bylaws \S4.3 no.\ 7 (member \(\to\) Team Lead \(\to\) Project Lead \(\to\) responsible Board member). No fixed Deputy Chair or CTO rule applies here, except where the CTO is in fact the Nearest Leader in that chain, or is acting as Deputy Chair for the CEO under \S12.3.

\vsection{7. Core Membership}
Receiving a genuine leadership responsibility through this contract makes the recipient a \textbf{Core Member (kjernemedlem)} (proposed New Bylaws \S5.1 no.\ 2); being asked to complete a task does not. This contract is used only for a genuine, standing leadership responsibility (e.g.\ Project Lead, Team Lead, or delegated Signature Right), never for a one-off task.

\vsection{8. General Assembly \(\to\) CEO}
Where the General Assembly itself is the granting party (e.g.\ the Mandate the newly elected CEO holds under \S14.1), the contract reads:

\begin{quote}
  ``The General Assembly, per the protocol of \meetingdate, 26H-ECL-GA-PROT1, gives \rule{3.5cm}{0.4pt} (CEO) the responsibility of \rule{4cm}{0.4pt}.''
\end{quote}

The CEO signs as \textbf{To}. The Protocol Signers do not personally become \textbf{From}: the responsibility originates from the General Assembly as a whole, and the Protocol ID is the record of that origin.

\vsection{9. Board record}
Once signed, the Board logs this contract using the Board Record, 26H-ECL-GA-RESP2. This contract and the Board Record cross-reference each other by Document ID.

\vsection{10. Signatures}
\begin{tabularx}{\textwidth}{|>{\bfseries\raggedright\arraybackslash}p{4cm}|X|}
  \hline
  From (name, role)        & \rule{0pt}{1.2cm} \\
  \hline
  From signature and date  & \rule{0pt}{1.2cm} \\
  \hline
  To (name, role)          & \rule{0pt}{1.2cm} \\
  \hline
  To signature and date    & \rule{0pt}{1.2cm} \\
  \hline
  Board Record Document ID & \rule{0pt}{1cm}   \\
  \hline
\end{tabularx}

\end{document}
